% ============================================================
%  CSD201 Library - Cheat Sheet
%  Biên dịch bằng XeLaTeX hoặc LuaLaTeX (Overleaf: chọn compiler XeLaTeX)
% ============================================================
\documentclass[10pt,a4paper]{article}

\usepackage[a4paper,top=1.5cm,bottom=1.7cm,left=1.5cm,right=1.5cm,
            headsep=0.3cm,footskip=0.8cm]{geometry}
\usepackage{fontspec}
\usepackage{xcolor}
\usepackage{sansmath}
\usepackage{tabularx}
\usepackage{array}
\usepackage{colortbl}
\usepackage{fancyhdr}
\usepackage{lastpage}
\usepackage{fancyvrb}
\usepackage[most]{tcolorbox}

% ---------- Phông chữ: cỡ chữ cố định, thống nhất ----------
\setmainfont{Latin Modern Sans}[Ligatures=TeX]
\setsansfont{Latin Modern Sans}
\setmonofont{Latin Modern Mono}[Scale=0.90]
\renewcommand{\familydefault}{\sfdefault}
\renewcommand{\normalsize}{\fontsize{9.5pt}{12.5pt}\selectfont}
\normalsize
\sansmath

\setlength{\parindent}{0pt}
\setlength{\parskip}{0pt}
\setlength{\tabcolsep}{5pt}
\renewcommand{\arraystretch}{1.22}
\pagestyle{fancy}
\fancyhf{}
\renewcommand{\headrulewidth}{0pt}
\renewcommand{\footrulewidth}{0.4pt}
\renewcommand{\footrule}{\hbox to\headwidth{\color{rule}\leaders\hrule height \footrulewidth\hfill}}
\fancyfoot[L]{\footnotesize\color{muted} CSD201 Library \,·\, Cheat Sheet}
\fancyfoot[R]{\footnotesize\color{muted} Trang \thepage/\pageref{LastPage}}

% ---------- Màu sắc ----------
\definecolor{ink}{HTML}{1F2933}
\definecolor{muted}{HTML}{6B7785}
\definecolor{rule}{HTML}{D5DBE1}
\definecolor{zebra}{HTML}{F5F7FA}
\definecolor{codebg}{HTML}{F3F5F8}
\definecolor{sig}{HTML}{14365D}
\definecolor{cList}{HTML}{1D5FA8}
\definecolor{cStack}{HTML}{0F7B6C}
\definecolor{cHash}{HTML}{B45309}
\definecolor{cRec}{HTML}{7C3AED}
\definecolor{cTree}{HTML}{2E7D32}
\definecolor{cSort}{HTML}{C2410C}
\definecolor{cGraph}{HTML}{B91C4C}
\color{ink}
\arrayrulecolor{rule}

% ---------- Lệnh trợ giúp ----------
% \fn{...}  : chữ đơn cách cho tên hàm
\newcommand{\fn}[1]{{\color{sig}\texttt{#1}}}
% \ret{...} : phần "trả về" / ghi chú, màu nhạt hơn
\newcommand{\ret}[1]{{\color{muted}$\rightarrow$\,#1}}
\newcommand{\code}[1]{\texttt{#1}}

% Tiêu đề module: \module{màu}{Tên}{mô tả ngắn}
\newcommand{\module}[3]{%
  \begin{tcolorbox}[enhanced,colback=#1!8,colframe=#1!8,boxrule=0pt,
     left=0pt,right=6pt,top=3pt,bottom=3pt,arc=2pt,
     borderline west={3.2pt}{0pt}{#1},
     before skip=0pt,after skip=5pt]
    \hspace{8pt}{\fontsize{12pt}{14pt}\selectfont\bfseries\color{#1!85!black}#2}%
    \hfill{\small\color{muted}#3}
  \end{tcolorbox}}

% Hộp ví dụ code
\DefineVerbatimEnvironment{Ex}{Verbatim}{fontsize=\fontsize{8.6pt}{11.4pt}\selectfont,commandchars=\\\{\}}
\newtcolorbox{example}[1][Ví dụ]{%
  enhanced,colback=codebg,colframe=rule,boxrule=0.4pt,arc=2pt,
  left=6pt,right=6pt,top=3pt,bottom=3pt,
  before skip=0pt,after skip=5pt,
  title={#1},fonttitle=\footnotesize\bfseries\color{muted},
  coltitle=muted,colbacktitle=codebg,attach boxed title to top left={yshift=-2.2mm,xshift=6pt},
  boxed title style={boxrule=0pt,frame hidden,left=2pt,right=2pt,top=0pt,bottom=0pt}}

% Bảng hàm: 2 cột (chữ ký hàm | cách dùng)
\newcolumntype{S}{>{\raggedright\arraybackslash\hangindent=1.3em\hangafter=1 }p{8.9cm}}
\newcolumntype{D}{>{\raggedright\arraybackslash}X}
\newenvironment{api}[1]{%
  \par\nointerlineskip\noindent
  \tabularx{\linewidth}{@{}S D@{}}
    \rowcolor{#1!14}\footnotesize\bfseries\color{#1!70!black}\;Hàm &
    \footnotesize\bfseries\color{#1!70!black}Cách dùng\\
    \hline
}{%
    \hline
  \endtabularx\par}

% Dòng nhóm trong bảng
\newcommand{\grp}[1]{%
  \multicolumn{2}{@{}l@{}}{\cellcolor{zebra}\footnotesize\bfseries\color{muted}\;#1}\\}
% Dòng hàm: \m{chữ ký}{mô tả}
\newcommand{\m}[2]{\;\fn{#1} & #2\\\hline}

\begin{document}

% ============================ TIÊU ĐỀ ============================
\begin{tcolorbox}[enhanced,colback=sig,colframe=sig,boxrule=0pt,arc=3pt,
   left=12pt,right=12pt,top=8pt,bottom=8pt,before skip=0pt,after skip=6pt]
  {\color{white}%
   {\fontsize{20pt}{24pt}\selectfont\bfseries CSD201 Library}\hfill
   {\large Cheat Sheet}\\[2pt]
   {\small\color{white!80!sig} Java \,·\, List \,·\, Stack/Queue \,·\, Hash Table \,·\,
     Recursion \,·\, BST \,·\, Sorting \,·\, Graph}}
\end{tcolorbox}

% ============================ 1. LIST ============================
\module{cList}{1 \,·\, MyList}{package \code{list} \,·\, danh sách liên kết đơn, lưu \code{int}}

\begin{api}{cList}
\grp{Khởi tạo \& thông tin}
\m{MyList()}{Tạo danh sách rỗng.}
\m{boolean isEmpty()}{Kiểm tra rỗng.}
\m{int size()}{Số node hiện có.}
\m{void traversal()}{In các giá trị từ head đến tail.}
\grp{Thêm}
\m{void addFirst(int x)}{Thêm \code{x} vào đầu danh sách.}
\m{void addLast(int x)}{Thêm \code{x} vào cuối danh sách.}
\m{boolean insertAt(int index, int x)}{Chèn tại \code{index} (từ \code{0} đến \code{size()}). \ret{\code{false} nếu sai vị trí}}
\m{boolean insertBefore(int targetValue, int newValue)}{Chèn \code{newValue} trước node đầu tiên có giá trị \code{targetValue}.}
\m{boolean insertAfter(int index, int x)}{Chèn \code{x} ngay sau node ở vị trí \code{index}.}
\grp{Truy cập \& cập nhật}
\m{Node get(int index)}{Lấy node tại \code{index}. \ret{\code{null} nếu sai vị trí}}
\m{Node search(int x)}{Tìm node đầu tiên có giá trị \code{x}. \ret{\code{null} nếu không thấy}}
\m{int getValue(int index)}{Lấy giá trị tại \code{index}. \ret{ném \code{IndexOutOfBoundsException}}}
\m{boolean updateAt(int index, int newValue)}{Đổi giá trị của node tại \code{index}.}
\grp{Xóa}
\m{boolean removeValue(int x)}{Xóa node đầu tiên có giá trị \code{x}.}
\m{boolean removeAt(int index)}{Xóa node tại \code{index}.}
\m{void remove(Node target)}{Xóa node \code{target} nếu thuộc danh sách.}
\m{int removeAllLessThan(int x)}{Xóa mọi node có giá trị $<$ \code{x}. \ret{số node đã xóa}}
\m{boolean removeAfter(int x)}{Xóa node \textbf{thứ hai} có giá trị $<$ \code{x}.}
\m{void clear()}{Xóa toàn bộ, đưa danh sách về rỗng.}
\grp{Min / Max \& sắp xếp}
\m{int min()}{Giá trị nhỏ nhất. \ret{ném \code{IllegalStateException} nếu rỗng}}
\m{int max()}{Giá trị lớn nhất. \ret{ném \code{IllegalStateException} nếu rỗng}}
\m{boolean swapMinMax()}{Đổi chỗ giá trị min đầu tiên và max đầu tiên.}
\m{void sort()}{Sắp xếp tăng dần (đổi giá trị, giữ liên kết). \ret{$O(n^2)$}}
\m{boolean sortRange(int startIndex, int endIndex)}{Sắp xếp đoạn \code{[startIndex, endIndex]} (gồm cả hai đầu).}
\m{void sortFirstThree()}{Sắp xếp 3 node đầu; ít hơn 3 node thì giữ nguyên.}
\m{void reverse()}{Đảo ngược bằng cách đổi liên kết. \ret{$O(n)$, bộ nhớ phụ $O(1)$}}
\m{boolean appendAll(MyList other)}{Nối bản sao giá trị của \code{other}. \ret{\code{false} nếu \code{null} / chính nó}}
\end{api}

\newpage
% ====================== 2. STACK & QUEUE ======================
\module{cStack}{2 \,·\, MyStack \& MyQueue}{package \code{stackAndQueue} \,·\, lưu \code{Object}}

\begin{api}{cStack}
\grp{MyStack \;(LIFO: vào sau, ra trước)}
\m{MyStack()}{Tạo stack rỗng.}
\m{boolean isEmpty()}{Kiểm tra rỗng.}
\m{void push(Object info)}{Đẩy \code{info} lên đỉnh stack.}
\m{Object pop()}{Lấy và xóa phần tử đỉnh. \ret{\code{null} nếu rỗng}}
\m{Object top()}{Xem phần tử đỉnh, không xóa. \ret{\code{null} nếu rỗng}}
\m{void clear()}{Xóa toàn bộ phần tử.}
\grp{MyQueue \;(FIFO: vào trước, ra trước)}
\m{MyQueue()}{Tạo queue rỗng.}
\m{boolean isEmpty()}{Kiểm tra rỗng.}
\m{void enqueue(Object info)}{Thêm \code{info} vào cuối hàng đợi.}
\m{Object dequeue()}{Lấy và xóa phần tử đầu. \ret{\code{null} nếu rỗng}}
\m{Object front()}{Xem phần tử đầu, không xóa. \ret{\code{null} nếu rỗng}}
\m{void clear()}{Xóa toàn bộ phần tử.}
\end{api}

\vspace{12pt}
% ============================ 3. HASH ============================
\module{cHash}{3 \,·\, MyHashTable}{package \code{hashTable} \,·\, nối chuỗi (chaining), $M=10$ bucket}

\begin{api}{cHash}
\m{MyHashTable()}{Tạo 10 bucket rỗng. \ret{hash: \code{floorMod(x, 10)}}}
\m{void add(int x)}{Thêm \code{x} vào bucket tương ứng. \ret{cho phép trùng}}
\m{void traversal()}{In các bucket không rỗng kèm chỉ số bucket.}
\m{int search(int x)}{Tìm \code{x}. \ret{trả về \code{x} nếu thấy, \code{-1} nếu không}}
\m{boolean remove(int x)}{Xóa \textbf{một} lần xuất hiện đầu tiên của \code{x}.}
\end{api}

\vspace{12pt}
% ========================== 4. RECURSION ==========================
\module{cRec}{4 \,·\, Recursion}{package \code{recursion} \,·\, toàn bộ là hàm \code{static}}

\begin{api}{cRec}
\m{static int factorial(int n)}{Tính $n!$. \ret{trả về \code{1} khi $n \le 1$}}
\m{static void printRecursivePattern(int n)}{In $n$ giảm dần rồi tăng dần. \ret{\code{n=3} in \code{321123}}}
\m{static String generatePattern(int n)}{Sinh chuỗi theo công thức $f(n)=f(n-3)+n+f(n-2)+n$.}
\m{static int calculateSum(int n)}{Tổng $1+2+\dots+n$. \ret{\code{0} nếu $n \le 0$}}
\m{static int findMin(int[] a, int n)}{Giá trị nhỏ nhất trong \code{n} phần tử đầu của \code{a}.}
\m{static boolean isPalindrome(char[] characters, int n)}{Kiểm tra \code{n} ký tự đầu có đối xứng không.}
\end{api}

\newpage
% ============================ 5. BST ============================
\module{cTree}{5 \,·\, BSTree}{package \code{tree} \,·\, cây nhị phân tìm kiếm, không nhận giá trị trùng}

\begin{api}{cTree}
\grp{Node \;(\code{info}, \code{left}, \code{right})}
\m{Node(int info)}{Tạo node không có con.}
\m{Node(int info, Node left, Node right)}{Tạo node kèm hai con.}
\grp{Cơ bản}
\m{BSTree()}{Tạo cây rỗng (\code{root = null}).}
\m{boolean isEmpty()}{Kiểm tra cây rỗng.}
\m{Node search(int key)}{Tìm node có giá trị \code{key}. \ret{\code{null} nếu không thấy}}
\m{void insert(int value)}{Thêm \code{value}; bỏ qua nếu đã tồn tại.}
\m{void visit(Node node)}{In giá trị của \code{node}.}
\m{int height(Node node)}{Chiều cao cây con. \ret{rỗng $=0$, chỉ có gốc $=1$}}
\grp{Duyệt cây \;(truyền vào \code{tree.root})}
\m{void preOrder(Node node)}{Gốc $\to$ Trái $\to$ Phải.}
\m{void inOrder(Node node)}{Trái $\to$ Gốc $\to$ Phải. \ret{BST in ra tăng dần}}
\m{void postOrder(Node node)}{Trái $\to$ Phải $\to$ Gốc.}
\m{void levelOrderTraversal(Node node)}{Duyệt theo từng tầng (dùng queue).}
\grp{Xóa, cân bằng \& xoay}
\m{void deleteByCopy(int value)}{Xóa; node 2 con được thay bằng \textbf{giá trị lớn nhất} bên trái.}
\m{void deleteByMerging(int value)}{Xóa; cây con trái thế chỗ, cây con phải gắn vào node phải nhất bên trái.}
\m{void balance()}{Cân bằng lại toàn bộ cây.}
\m{Node rotateLeft(Node node)}{Xoay trái. \ret{gốc mới của cây con}}
\m{Node rotateRight(Node node)}{Xoay phải. \ret{gốc mới của cây con}}
\end{api}

\vspace{12pt}
% ============================ 6. SORTING ============================
\module{cSort}{6 \,·\, Sorting}{package \code{sorting} \,·\, tất cả sắp xếp tăng dần, hàm của đối tượng}

\begin{api}{cSort}
\m{void mergeSort(int[] array)}{Merge Sort. \ret{$O(n\log n)$, bộ nhớ phụ $O(n)$}}
\m{void quickSort(int[] array)}{Quick Sort. \ret{TB $O(n\log n)$, xấu nhất $O(n^2)$}}
\m{void selectionSort(int[] array)}{Selection Sort. \ret{$O(n^2)$}}
\m{void insertionSort(int[] array)}{Insertion Sort. \ret{$O(n^2)$}}
\m{void bubbleSort(int[] array)}{Bubble Sort, dừng sớm nếu đã sắp xếp. \ret{$O(n^2)$}}
\m{void swap(int[] array, int firstIndex, int secondIndex)}{Đổi chỗ hai phần tử trong mảng.}
\m{void saveToFile(int[] array, String fileName)}{Ghi mảng ra file, mỗi phần tử một dòng. \ret{ném \code{IOException}}}
\end{api}

\newpage
% ============================ 7. GRAPH ============================
\module{cGraph}{7 \,·\, Graph}{package \code{graph} \,·\, ma trận kề \,·\, đỉnh đánh số từ 0 (0 $\to$ A, 1 $\to$ B, \dots)}

\begin{api}{cGraph}
\grp{Khởi tạo \;(\code{Graph.NO\_EDGE} = \code{Integer.MAX\_VALUE / 4})}
\m{Graph(int vertexCount)}{Đồ thị vô hướng, không trọng số.}
\m{Graph(int vertexCount, boolean weighted, boolean directed)}{Đồ thị tùy chọn có trọng số / có hướng.}
\m{Graph(int[][] adjacencyMatrix)}{Từ ma trận kề vô hướng; \code{0} = không có cạnh.}
\m{Graph(int[][] adjacencyMatrix, boolean weighted, boolean directed, int noEdgeValue)}{Từ ma trận kề, tự chỉ định giá trị biểu diễn \textbf{không có cạnh}.}
\grp{Thông tin}
\m{int getVertexCount()}{Số đỉnh.}
\m{boolean isWeighted()}{Đồ thị có trọng số không.}
\m{boolean isDirected()}{Đồ thị có hướng không.}
\m{void visit(int vertex)}{In nhãn đỉnh (A, B, \dots), không xuống dòng.}
\m{void printMatrix()}{In ma trận kề. \ret{\code{-} = không có cạnh, đường chéo = 0}}
\grp{Thêm / xóa cạnh}
\m{void addEdge(int from, int to)}{Thêm cạnh với trọng số mặc định \code{1}.}
\m{void addEdge(int from, int to, int weight)}{Thêm cạnh có trọng số. \ret{vô hướng: lưu hai chiều, bỏ self-loop}}
\m{void removeEdge(int from, int to)}{Xóa cạnh (vô hướng: xóa cả hai chiều).}
\grp{Duyệt \& tìm đường}
\m{void breadthFirstTraversal(int start)}{BFS (dùng queue); đồ thị không liên thông thì duyệt tiếp đỉnh chưa thăm.}
\m{void depthFirstTraversal(int start)}{DFS (đệ quy); cũng duyệt các thành phần còn lại.}
\m{void printConnectedComponents()}{In các thành phần liên thông. \ret{chỉ đồ thị vô hướng}}
\m{void printPath(int start, int destination)}{In \textbf{một} đường đi bằng DFS. \ret{không đảm bảo ngắn nhất}}
\m{void findShortestPath(int start, int destination)}{Dijkstra: in khoảng cách và đường đi. \ret{trọng số $\ge 0$}}
\m{void findEulerCircuit()}{Tìm chu trình Euler. \ret{vô hướng, mọi đỉnh bậc chẵn, liên thông; không cạnh song song/self-loop}}
\end{api}

\end{document}
